\begin{center}\textit{A new algorithm for deciding whether a given positive integer is a quadratic residue or a nonresidue of a given positive prime number.}\end{center}

\Article{1}

Before presenting a new solution to this problem, we shall briefly recall the solution given in the \textit{Disquisitiones Arithmeticae}. It can be carried out quite readily with the help of the fundamental theorem and the following known theorems:

I. The relation of the number $a$ to the number $b$ (whether the former is a quadratic residue or a nonresidue of the latter) is the same as that of the number $c$ to $b$, if $a\equiv c\ (\mathrm{mod.}\ b)$.

II. If $a$ is the product of the factors $\alpha,\beta,\gamma,\delta$ etc., and $b$ is prime, the relation of $a$ to $b$ depends on the relations of these factors to $b$ as follows: $a$ is a quadratic residue of $b$ or a nonresidue according as an even or an odd number of those factors are nonresidues of $b$. Thus a factor that is a square can be disregarded entirely in this examination; if a factor is an integer raised to an odd power, that integer itself can take its place.

III. The number 2 is a quadratic residue of every prime of the form $8m+1$ or $8m+7$, and a nonresidue of every prime of the form $8m+3$ or $8m+5$.

\SourceLeaf{70}%
Suppose, then, that a number $a$ is given whose relation to a prime $b$ is sought. If $a$ is greater than $b$, first replace it by its least positive residue modulo $b$. Once this residue has been resolved into prime factors, theorem II reduces the question to finding the relation of each factor to $b$. The relation of the factor $2$ (if it occurs once, three times, five times, etc.) is known from theorem III; by the fundamental theorem, the relation of each remaining factor depends on the relation of $b$ itself to that factor. Thus, in place of the single relation of the given number to the prime $b$, we must now investigate several relations of the number $b$ to other odd primes smaller than $b$. These problems are reduced in the same way to smaller moduli, and these successive reductions must evidently come to an end.

\Article{2}

To illustrate this solution, suppose that the relation of the number $103$ to $379$ is sought. Since $103$ is already less than $379$ and is itself prime, the fundamental theorem applies immediately: it tells us that the required relation is opposite to that of the number $379$ to $103$. This in turn is the same as the relation of the number $70$ to $103$, which itself depends on the relations of $2,\ 5,\ 7$ to $103$. The first of these is known from theorem III. By the fundamental theorem, the second depends on the relation of the number $103$ to $5$, which by theorem I is the same as that of $3$ to $5$; this in turn depends, by the fundamental theorem, on the relation of $5$ to $3$, which by theorem I is the same as that of $2$ to $3$, known from theorem III. Likewise, by the fundamental theorem, the relation of $7$ to $103$ depends on that of $103$ to $7$, which by theorem I is the same as that of $5$ to $7$; this in turn depends, by the fundamental theorem, on the relation of $7$ to $5$, which by theorem I is the same as that of $2$ to $5$, known from theorem III. If we now wish to turn this analysis into a synthesis, the decision is reduced to fourteen steps. We give them all here, so that the greater elegance of the new solution may stand out more clearly.

{\raggedright\hangindent=2.2em\hangafter=1\noindent 1.\quad The number $2$ is a quadratic residue of the number $103$ \mbox{(theorem III)}.\par}

{\raggedright\hangindent=2.2em\hangafter=1\noindent 2.\quad The number $2$ is a quadratic nonresidue of the number $3$ \mbox{(theorem III)}.\par}

{\raggedright\hangindent=2.2em\hangafter=1\noindent 3.\quad The number $5$ is a quadratic nonresidue of the number $3$ \mbox{(from I and 2)}.\par}

{\raggedright\hangindent=2.2em\hangafter=1\noindent 4.\quad The number $3$ is a quadratic nonresidue of the number $5$ \mbox{(the fundamental theorem and 3)}.\par}

{\raggedright\hangindent=2.2em\hangafter=1\noindent 5.\quad The number $103$ is a quadratic nonresidue of the number $5$ \mbox{(I and 4)}.\par}

\SourceLeaf{71}%
{\raggedright\hangindent=2.2em\hangafter=1\noindent 6.\quad The number $5$ is a quadratic nonresidue of the number $103$ \mbox{(the fundamental theorem and 5)}.\par}

{\raggedright\hangindent=2.2em\hangafter=1\noindent 7.\quad The number $2$ is a quadratic nonresidue of the number $5$ \mbox{(theorem III)}.\par}

{\raggedright\hangindent=2.2em\hangafter=1\noindent 8.\quad The number $7$ is a quadratic nonresidue of the number $5$ \mbox{(I and 7)}.\par}

{\raggedright\hangindent=2.2em\hangafter=1\noindent 9.\quad The number $5$ is a quadratic nonresidue of the number $7$ \mbox{(the fundamental theorem and 8)}.\par}

{\raggedright\hangindent=2.2em\hangafter=1\noindent 10.\quad The number $103$ is a quadratic nonresidue of the number $7$ \mbox{(I and 9)}.\par}

{\raggedright\hangindent=2.2em\hangafter=1\noindent 11.\quad The number $7$ is a quadratic residue of the number $103$ \mbox{(the fundamental theorem and 10)}.\par}

{\raggedright\hangindent=2.2em\hangafter=1\noindent 12.\quad The number $70$ is a quadratic nonresidue of the number $103$ \mbox{(II, 1, 6, 11)}.\par}

{\raggedright\hangindent=2.2em\hangafter=1\noindent 13.\quad The number $379$ is a quadratic nonresidue of the number $103$ \mbox{(I and 12)}.\par}

{\raggedright\hangindent=2.2em\hangafter=1\noindent 14.\quad The number $103$ is a quadratic residue of the number $379$ \mbox{(the fundamental theorem and 13)}.\par}

For brevity, in what follows we shall use the notation introduced in \textit{Comment. Gotting. Vol. XVI}. By $[x]$ we shall mean the quantity $x$ itself when $x$ is an integer, and the greatest integer below $x$ when $x$ is not an integer. Thus $x-[x]$ is always nonnegative and less than unity.

\Article{3}

\textsc{Problem.}\quad \textit{Let $a,\ b$ be positive integers coprime to one another, and put $[\frac12a]=a'$. Find the sum}

\[
\gbr{\frac ba}+\gbr{\frac{2b}a}+\gbr{\frac{3b}a}+\gbr{\frac{4b}a}+\text{etc.}+\gbr{\frac{a'b}a}
\]

\Needspace{9\baselineskip}
\textsc{Solution.}\quad For brevity, let us denote such a sum by $\varphi(a,b)$, so that we also have

\[
\varphi(b,a)=\gbr{\frac ab}+\gbr{\frac{2a}b}+\gbr{\frac{3a}b}+\text{etc.}+\gbr{\frac{b'a}b}
\]

if we put $[\frac12b]=b'$. In the third proof of the fundamental theorem it was shown that, when $a$ and $b$ are odd,

\[
\varphi(a,b)+\varphi(b,a)=a'b'
\]

By following the same method, the truth of this proposition extends easily to the case in which one of the numbers $a,\ b$ is odd, as we already indicated there. Divide $a$ by $b$, as in the method for finding the greatest common divisor of two integers, and let $\beta$ be the quotient and $c$ the remainder; then divide $b$ by $c$, and so on, obtaining the equations

\SourceLeaf{72}%
\[
\begin{aligned}
a&=\beta b+c\\
b&=\gamma c+d\\
c&=\delta d+e\\
d&=\varepsilon e+f\ \text{etc.}
\end{aligned}
\]

In this way the continually decreasing sequence of numbers $b,c,d,e,f$ etc. will eventually reach unity, since by hypothesis $a$ and $b$ are coprime. The last equation is therefore

\[
k=\lambda l+1
\]

Since evidently

\[
\begin{aligned}
\gbr{\frac ab}&=\gbr{\beta+\frac cb}=\beta+\gbr{\frac cb}\\
\gbr{\frac{2a}b}&=\gbr{2\beta+\frac{2c}b}=2\beta+\gbr{\frac{2c}b}\\
\gbr{\frac{3a}b}&=\gbr{3\beta+\frac{3c}b}=3\beta+\gbr{\frac{3c}b}
\end{aligned}
\]

etc.. we have

\[
\varphi(b,a)=\varphi(b,c)+\tfrac12\beta(b'b'+b')
\]

and consequently

\[
\varphi(a,b)=a'b'-\tfrac12\beta(b'b'+b')-\varphi(b,c)
\]

\Needspace{10\baselineskip}
By similar reasoning, on putting $[\frac12c]=c',\ [\frac12d]=d'\ [\frac12e]=e'$ etc., we obtain

\[
\begin{aligned}
\varphi(b,c)&=b'c'-\tfrac12\gamma(c'c'+c')-\varphi(c,d)\\
\varphi(c,d)&=c'd'-\tfrac12\delta(d'd'+d')-\varphi(d,e)\mathord{.}\\
\varphi(d,e)&=d'e'-\tfrac12\varepsilon(e'e'+e')-\varphi(e,f)
\end{aligned}
\]

etc., as far as

\[
\varphi(k,l)=k'l'-\tfrac12\lambda(l'l'+l')-\varphi(l,1)
\]

Since evidently $\varphi(l,1)=0$, we obtain from this the formula

\[
\begin{aligned}
\varphi(a,b)&=a'b'-b'c'+c'd'-d'e'+\text{etc.}\pm k'l'\\
&\quad-\tfrac12\beta(b'b'+b')+\tfrac12\gamma(c'c'+c')-\tfrac12\delta(d'd'+d')\\
&\quad+\tfrac12\varepsilon(e'e'+e')-\text{etc.}\mp\tfrac12\lambda(l'l'+l')
\end{aligned}
\]

\Article{4}

It now follows easily from the third proof that the relation of the number $b$ to $a$, whenever $a$ is prime, can be read directly from the value of the sum\SourceLeaf{73} $\varphi(a,2b)$. According as this sum is even or odd, $b$ is a quadratic residue of $a$ or a nonresidue. The sum $\varphi(a,b)$ itself can also be used for the same purpose, provided that we distinguish the case in which $b$ is odd from that in which it is even. Namely,

I.\quad When $b$ is odd, $b$ is a quadratic residue or a nonresidue of $a$ according as $\varphi(a,b)$ is even or odd.

II.\quad When $b$ is even, the same rule holds if, in addition, $a$ is of the form $8n+1$ or $8n+7$; but if, for an even value of $b$, the modulus $a$ is of the form $8n+3$ or $8n+5$, the opposite rule must be applied: $b$ is a quadratic residue of $a$ if $\varphi(a,b)$ is odd, and a nonresidue if $\varphi(a,b)$ is even.

All this follows very easily from Article 4 of the third proof.

\newpage
\Article{5}

\textit{Example.} To find the relation of the number $103$ to the prime $379$, we have, for calculating the sum $\varphi(379,103)$,

\[
\begin{array}{r|r|l}
a=379&a'=189&\\
b=103&b'=51&\beta=3\\
c=70&c'=35&\gamma=1\\
d=33&d'=16&\delta=2\\
e=4&e'=2&\varepsilon=8
\end{array}
\]

and hence

\[
\begin{aligned}
\varphi(379,103)&=9639-1785+560-32-3978+630-272+24\\
&=4786
\end{aligned}
\]

so $103$ is a quadratic residue of the number $379$. If we prefer to use the sum $(379,206)$ for the same purpose, we have the following scheme:

\[
\begin{array}{r|r|l}
379&189&\\
206&103&1\\
173&86&1\\
33&16&5\\
8&4&4
\end{array}
\]

from which we obtain

\SourceLeaf{74}%
\[
\begin{aligned}
\varphi(379,206)&=19467-8858+1376-64-5356+3741-680+40\\
&=9666
\end{aligned}
\]

Thus $103$ is a quadratic residue of the number $379$.

\Needspace{28\baselineskip}
\Article{6}

To determine the relation of the number $b$ to $a$, it is not necessary to calculate the individual terms of the sum $\varphi(a,b)$; it is enough to know how many of them are odd. Our rule can therefore also be stated as follows:

As above, take $a=\beta b+c,\ b=\gamma c+d,\ c=\delta d+e$ etc., until the sequence of numbers $a,\ b,\ c,\ d,\ e$ etc. reaches unity. Put $[\frac12a]=a',\ [\frac12b]=b',\ [\frac12c]=c'$ etc., and let $\mu$ count the odd numbers in the sequence $a',b',c'$ etc. that are immediately followed by an odd number. Further, let $\nu$ count the odd numbers in the sequence $\beta,\gamma,\delta$ etc. whose corresponding numbers in the sequence $b',c',d'$ etc., respectively, have the form $4n+1$ or $4n+2$. Then $b$ is a quadratic residue or a nonresidue of $a$ according as $\mu+\nu$ is even or odd, with the sole exception that, when $b$ is even and $a$ is of the form $8n+3$ or $8n+5$, the opposite rule holds.

In our example the sequence $a',b',c',d',e'$ has two instances in which an odd entry is immediately followed by an odd entry, so $\mu=2$; the sequence $\beta',\gamma',\delta',\varepsilon'$ does contain two odd numbers, but their corresponding numbers in the sequence $b',c',d',e'$ have the form $4n+3$, so $\nu=0$. Thus $\mu+\nu$ is even, and consequently $103$ is a quadratic residue of the number $379$.

\BookRule{42mm}
